% SGA 3, Expose V, Loop 1 English reconstruction.
% Scope: component-local pp. 36--38, continuation of Theorem 10.1.2
% through the statement of Theorem 10.4.1.
% Authority: Polo--Gille Exp5-13oct24.pdf, combined-reader pp. 298--300,
% re-edited headers 286--288 (original running markers 289--291).
% Jacob Reinhold's CC BY 4.0 Markdown is comparison/drafting material only.

\textit{分解为}
\SGAThreeDiagram[0.42\textwidth]
  {figures/exp5/p286_theorem1012_factorization.png}
  {定理~\hyperref[thm:sga3-V-10.1.2]{10.1.2}
   中轨道态射的分解.}
  {fig:sga3-V-10.1.2-factorization}
  {第 V 篇, 组件局部第~36 页 / 合订阅读版第~298 页 /
   修订版页眉第~286 页}
\textit{其中 \(p\) 为典范投影, 它忠实
平坦且局部有限呈示, \(i\) 为单态射.}

\medskip
\noindent\textit{证明.}---
态射 \(f\) 使 \(G\) 成为一个 \(X\)-概形. 由
\(\xi\) 的稳定化子的定义, 态射
\phantomsection\label{eq:sga3-V-10.1.2-stabilizer-isomorphism}%
\[
  G\times_S H\longrightarrow G\times_XG,
  \qquad (g,h)\longmapsto(g,gh)
\]
是同构. 由于 \(H\) 在 \(S\) 上平坦, 概形
\(G\times_SH\) 在 \(G\) 上平坦; 因而第一投影
\[
  p_1:G\times_XG\longrightarrow G
\]
平坦.%
\footnote{\label{source-note:sga3-V-10.1.2-fiber-product-base}源 PDF 在此投影中印作 \(G\times_SG\).
\hyperref[eq:sga3-V-10.1.2-stabilizer-isomorphism]{前面的同构公式}及定理~\hyperref[thm:sga3-V-10.1.1]{10.1.1}
的应用要求纤维积为 \(G\times_XG\), 此处采用这一读法.}
此外, 若 \(X\) 在 \(S\) 上局部有限型, 则 \(f\)
局部有限呈示 (EGA~IV\(_1\), 1.4.3(v)); 否则,
已假设 \(S\) 局部 Noether. 因此只需对 \(f\) 应用
定理~\hyperref[thm:sga3-V-10.1.1]{10.1.1}. 只剩下
验证 \(G/H\) 在 \(S\) 上局部有限呈示,
而这立即由命题~\hyperref[sga3:v:proposition:9:1]{9.1} 得到.

\medskip
\phantomsection\label{cor:sga3-V-10.1.3}%
\noindent\textbf{推论 10.1.3.}---
\textit{设 \(S\) 为概形, \(u:G\to H\) 为
\(S\)-群概形的态射. 假设 \(G\) 在 \(S\) 上局部有限
呈示, 且 \(H\) 在 \(S\) 上局部有限
型或 \(S\) 局部 Noether.}

\textit{若 \(K=\ker(u)\) 在 \(S\) 上平坦, 则商群
\(G/K\) 可由一个在 \(S\) 上局部有限呈示的
\(S\)-群概形表出, 且 \(u\) 分解为}
\SGAThreeDiagram[0.43\textwidth]
  {figures/exp5/p286_corollary1013_factorization.png}
  {\(u\) 经由 \(G/K\) 的分解.}
  {fig:sga3-V-10.1.3-factorization}
  {第 V 篇, 组件局部第~36 页 / 合订阅读版第~298 页 /
   修订版页眉第~286 页}
\textit{其中 \(p\) 为典范投影, \(i\) 为
单态射.}

\medskip
\noindent\textit{证明.}---
取 \(X=H\), 并取 \(H\) 的单位截面为 \(\xi\),
应用定理~\hyperref[thm:sga3-V-10.1.2]{10.1.2}.

\section*{10.2. 对角线的稳定化子}
\phantomsection\label{sec:sga3-V-10.2}
\addcontentsline{toc}{section}{\SGARefLink{sec:sga3-V-10.2}{10.2}. 对角线的稳定化子}

设 \(S\) 为 Noether 概形, \(X\) 为有限
型 \(S\)-概形, \(G\) 为有限型平坦 \(S\)-群概形,
左作用于 \(X\). 因此有一个 \(S\)-作用
\(d_0:G\times_SX\to X\). 设 \(d_1:G\times_SX\to X\) 为到第二个
因子的投影. 如 \hyperref[par:sga3-V-2-a]{\S2(a)} 中一样, 有群胚
\SGAThreeDiagram[0.67\textwidth]
  {figures/exp5/p286_section102_action_groupoid.png}
  {与 \((X,G)\) 相伴的作用群胚.}
  {fig:sga3-V-10.2-action-groupoid}
  {第 V 篇, 组件局部第~36 页 / 合订阅读版第~298 页 /
   修订版页眉第~286 页}
其余核若存在, 记为 \(G\backslash X\).

\medskip
\phantomsection\label{def:sga3-V-10.2.1}%
\noindent\textbf{定义 10.2.1.}---
\textit{用 \(F\subset G\times_SX\) 表示\emph{对角线截面的稳定化子};
即 \(F\) 是由下列笛卡尔方块定义的 \(X\)-概形:}
\SGAThreeDiagram[0.47\textwidth]
  {figures/exp5/p287_definition1021_cartesian_square.png}
  {对角线截面的稳定化子 \(F\).}
  {fig:sga3-V-10.2.1-stabilizer}
  {第 V 篇, 组件局部第~37 页 / 合订阅读版第~299 页 /
   修订版页眉第~287 页}

于是 \(F\) 是 \(G\times_SX\) 的 \(X\)-子群概形. 由于
\(G\times_SX\) 在 Noether 概形 \(S\) 上有限型,
它是 Noether 概形; 因此 \(F\) 在 \(S\)
及 \(X\) 上都有限型 (EGA~I, 6.3.5 及 6.3.6). 若 \(X\to S\) 分离,
则 \(F\) 是 \(G\times_SX\) 的闭 \(X\)-子群概形.

回顾, 当
\[
  (d_0,d_1):G\times_SX\longrightarrow X\times_SX
\]
为单态射时, 称 \(G\) 在 \(X\) 上\emph{自由作用}
(第~III 篇, 3.2.1). 等价地, \(F\) 是
\(X\) 上的平凡群概形.

\section*{10.3. \(F\) 在 \(X\) 上拟有限的情形}
\phantomsection\label{sec:sga3-V-10.3}
\addcontentsline{toc}{section}{\SGARefLink{sec:sga3-V-10.3}{10.3}. F 在 X 上拟有限的情形}

由于 \(F\) 在 \(X\) 上有限型, 它在
\(X\) 上拟有限当且仅当 \(X\) 的几何点的稳定化子
有限.

\medskip
\phantomsection\label{thm:sga3-V-10.3.1}%
\noindent\textbf{定理 10.3.1.}---%
\footnote{\label{note:sga3-V-49}编者注: 这里也存在一个
满足\hyperref[thm:sga3-V-10.3.1]{此定理}结论的最大开子集 \(U\subset X\);
见编者\hyperref[note:sga3-V-44]{注~44}.}
\textit{在 \hyperref[sec:sga3-V-10.2]{10.2} 的假设下,
假设 \(F\) 在 \(X\) 上拟有限. 则存在一个稠密的
\(G\)-饱和开子集 \(U\subset X\), 具有下列
性质:}
\begin{enumerate}
\item[(i)]\phantomsection\label{thm:sga3-V-10.3.1-i}%
  \textit{在 \((\mathbf{Sch}/S)\) 中, 余核
  \(V=G\backslash U\) 存在; 此外, \(V\) 是
  赋环空间范畴中的商;}
\item[(ii)]\phantomsection\label{thm:sga3-V-10.3.1-ii}%
  \textit{\(p:U\to V\) 满射、为开映射且有限
  呈示;}
\item[(iii)]\phantomsection\label{thm:sga3-V-10.3.1-iii}%
  \textit{\(V\) 在 \(S\) 上有限呈示;}
\item[(iv)]\phantomsection\label{thm:sga3-V-10.3.1-iv}%
  \textit{态射
  \[
    G\times_SU\longrightarrow U\times_VU,
    \qquad (g,x)\longmapsto(gx,x),
  \]
  满射;}
\item[(v)]\phantomsection\label{thm:sga3-V-10.3.1-v}%
  \textit{若此外 \(G\) 在 \(X\) 上自由作用, 则
  \(U\to V\) 是左 \(G\)-挠子 (torsor), 在 fppf
  拓扑下局部平凡. 特别地, \(U\to V\) 忠实平坦.}%
  \footnote{\label{note:sga3-V-50}编者注: 若还假设
  \(G\) 是约化 \(S\)-群概形, 且其在 \(X\) 上的
  自由作用可线性化, 则已知 \(G\backslash X\)
  可表, 且 \(X\to G\backslash X\) 是左
  \(G\)-挠子. 这由 Raynaud 和 Seshadri 的结果得到,
  见 \(\hyperref[bib:sga3-V-CTS79]{[\mathrm{CTS79}]}\), 命题~6.11.}
\end{enumerate}

\noindent\textit{证明.}---
由假设, 态射
\[
  G\times_SX\longrightarrow X\times_SX,
  \qquad (g,x)\longmapsto(gx,x),
\]
拟有限. 因此可将定理~\hyperref[thm:sga3-V-8.1]{8.1}
应用于 \((X,G)\) 定义的群胚. 从而存在一个稠密
饱和开子集 \(U\subset X\), 使得商
\(G\backslash U\) 存在且具有性质 \hyperref[thm:sga3-V-10.3.1-i]{(i)}, \hyperref[thm:sga3-V-10.3.1-ii]{(ii)} 和 \hyperref[thm:sga3-V-10.3.1-iii]{(iii)};
\hyperref[thm:sga3-V-8.1]{同一定理}还给出 \hyperref[thm:sga3-V-10.3.1-iv]{(iv)}.

为证明 \hyperref[thm:sga3-V-10.3.1-v]{(v)}, 回顾 \(G\) 在 \(X\) 上自由作用当且仅当
\((d_0,d_1)\) 为等价对 (第~III 篇, 3.2.1).%
\footnote{\label{source-note:sga3-V-10.3.1-proof-part-v}源 PDF 此处写作\emph{为证明 (iv)}.
所引用的定理~\hyperref[thm:sga3-V-8.1-iv]{8.1(iv)} 及关于挠子的结论证明的是
当前编号中的 \hyperref[thm:sga3-V-10.3.1-v]{(v) 部分}; \hyperref[thm:sga3-V-10.3.1-iv]{(iv) 部分}由
定理~\hyperref[thm:sga3-V-8.1-iii]{8.1(iii)} 得到.}
此时定理~\hyperref[thm:sga3-V-8.1-iv]{8.1(iv)} 表明
\[
  G\times_SU\longrightarrow U\times_VU
\]
是同构, 且 \(p\) 忠实平坦且有限
呈示. 因此 \(U\) 是 \(V\) 上的 \(G\)-挠子,
在 fppf 拓扑下局部平凡.

\section*{10.4. \(F\) 在 \(X\) 上平坦的情形}
\phantomsection\label{sec:sga3-V-10.4}
\addcontentsline{toc}{section}{\SGARefLink{sec:sga3-V-10.4}{10.4}. F 在 X 上平坦的情形}

记
\[
  d=(d_0,d_1):G\times_SX\longrightarrow X\times_SX,
  \qquad d(g,x)=(gx,x).
\]
回顾, 与 \((X,G)\) 相伴的等价关系的层论意义下的图
\(\widetilde\Gamma\), 是 \(X\times_SX\) 中由
\((d_0,d_1)\) 的像给出的 fppf
\(S\)-子层. 它是与图函子
\[
  T\longmapsto
  \Gamma(T)=\bigl\{(x_0,x_1)\in X(T)\times X(T)
  \mathrel{|} x_0\in G(T)x_1\bigr\}.
\]
相伴的 fppf 层.
令 \(G_X=G\times_SX\). 对每个 \(S\)-概形 \(T\), 有
满射
\[
  G_X(T)\longrightarrow\Gamma(T),
  \qquad (g,x)\longmapsto(gx,x),
\]
它诱导双射%
\phantomsection\label{eq:sga3-V-10.4-phiT-bijection}%
\[
  \phi(T):G_X(T)/F(T)\xrightarrow{\ \sim\ }\Gamma(T).
\]
事实上, 若 \((g,x),(g',x')\in G_X(T)\) 有相同的像, 则
\(x'=x\) 且 \(g^{-1}g'x=x\), 从而
\((g^{-1}g',x)\in F(T)\), 而 \((g,x)\) 和 \((g',x)\) 在
\(G_X(T)/F(T)\) 中有相同的像.%
\footnote{\label{source-note:sga3-V-10.4-ft-typing}源 PDF 印作
\((g^{-1}g'x,x)\in F(T)\). 由于
\(F(T)\subseteq G(T)\times X(T)\), 第一个坐标必须是
群元素 \(g^{-1}g'\); 英文版采用这一类型正确的
读法. 同一句中, 原文曾将分母写成 \(F_X(T)\);
\hyperref[eq:sga3-V-10.4-phiT-bijection]{上面的双射公式}与
\hyperref[eq:sga3-V-10.4-quotient-presheaf]{下面的预层}均使用已定义的稳定化子 \(F(T)\),
此处始终采用这一记号.}

由定义 (见 \SGARefLink{sga3:iv:theorem:4:4:1:item:ii:01}{第~IV 篇, 4.4.1(ii)}, 或 5.2.1 的证明),
商层 \(G_X/F\) 是与预层%
\phantomsection\label{eq:sga3-V-10.4-quotient-presheaf}%
\[
  T\longmapsto G_X(T)/F(T)\simeq\Gamma(T).
\]
相伴的 fppf 层.
因此得到层的同构
\[
  \phi:G_X/F\xrightarrow{\ \sim\ }\widetilde\Gamma.
\]

\medskip
\phantomsection\label{thm:sga3-V-10.4.1}%
\noindent\textbf{定理 10.4.1.}---%
\footnote{\label{note:sga3-V-51}编者注: 这是
\(\hyperref[bib:sga3-V-Ray67b]{[\mathrm{Ray67b}]}\) 的定理~3 的 (2) 部分. 该文还概述了
定理~\hyperref[thm:sga3-V-10.1.1]{10.1.1} 的另一个证明.}
\textit{在 \hyperref[sec:sga3-V-10.2]{10.2} 的假设下,
下列断言成立:}
\begin{enumerate}
\item[\textit{a)}]\phantomsection\label{thm:sga3-V-10.4.1-a}%
  \textit{层 \(\widetilde\Gamma\) 可表当且仅当
  \(F\) 在 \(X\) 上平坦;}
\item[\textit{b)}]\phantomsection\label{thm:sga3-V-10.4.1-b}%
  \textit{假设 \(F\) 在 \(X\) 上平坦. 则由 \(d_1\) 和 \(d_0\)
  诱导的态射}
  \SGAThreeDiagram[0.37\textwidth]
    {figures/exp5/p288_theorem1041_double_arrow.png}
    {由 \(d_1\) 和 \(d_0\) 诱导的两个映射.}
    {fig:sga3-V-10.4.1-double-arrow}
    {第 V 篇, 组件局部第~38 页 / 合订阅读版第~300 页 /
     修订版页眉第~288 页}
  \textit{忠实平坦且有限呈示.}
\end{enumerate}

% 翻译统计: 源行数 263; 输出行数 264 (含本统计行); 行内/行间数学片段 156; 翻译占位 0. 底本疑点另报主线程.
